\section{\uppercase{Introduction}}

\noindent Many agents that move through buildings must know where they are. A pedestrian
following turn-by-turn directions on a smartphone, a mobile robot running
warehouse logistics or facility inspection, and a tracked asset reporting its
location all depend on a position estimate. Outdoors that estimate comes from
GPS, but satellite signals scarcely penetrate walls. Consequently, an indoor device must
infer its position from the sensors it already carries. The inference is
error prone because no effective global reference is available indoors: the
measurements are noisy, and the environment shifts as people and objects move.
Accuracy is bounded further by the sensing hardware and by the site survey that
calibrates it. An estimate that holds under one set of conditions can degrade
sharply once any of them changes, so reliable indoor self-localization remains an
open problem.

Two sensors are carried by nearly all of these devices: a WiFi radio and an
inertial measurement unit (IMU). A WiFi scan gives a sparse absolute reference to
known access points. The IMU reports dense relative motion at a significantly higher rate.
The two streams are complementary, yet heterogeneous and sampled asynchronously.
No common clock aligns their arrivals, and either stream may become unavailable or arrive
stale during operation. Fusing these imperfect, asynchronous signals into one
position estimate is the core problem this work targets.

In this paper, we propose a continuous-time transformer (Figure~\ref{fig:pipeline})
that fuses asynchronous WiFi and IMU observations into a single $(x, y)$ estimate
within one model that remains accurate when a sensor is missing or stale. The
approach is not tied to a particular platform. It needs only WiFi scans and
inertial windows, the same inputs a phone, a robot, or a wearable already produces,
so it applies to any device that carries these two sensors. We validate it on real
cross-session data.

The remainder of this paper is organized as follows: Section~\ref{sec:related}
reviews related work, Section~\ref{sec:problem} states the problem,
Section~\ref{sec:method} develops the method, Section~\ref{sec:experiments}
describes the experimental setup, Section~\ref{sec:results} reports the results,
and Section~\ref{sec:conclusion} concludes.
